\chapter{Value function approximation}
\label{ch:vfa}

When the states are too many for a table, the value function is
represented by a parametrised function $\hat v(s; \mathbf{w})$ with far
fewer parameters than states. This chapter studies the simplest and best
understood case, \emph{linear} approximation, where the theory is complete:
it shows what TD learning converges to, how far this is from the best
possible approximation, and why off-policy learning can diverge. Deep
networks (\cref{ch:dqn}) have no such guarantees, but the phenomena
identified here explain much of their behaviour.

\section{The prediction objective}

We approximate the value $v_\pi$ of a fixed policy. Errors in different
states cannot all be made small, so we weight them by a distribution $d$
over states, usually the frequency with which the policy visits them.

\begin{definition}{Linear approximation and weighted norm}{linear-fa}
  A \emph{feature map} $\mathbf{x} : \States \to \R^m$ assigns to each state
  a vector of $m$ features; the approximation is $\hat v(s; \mathbf{w}) =
  \mathbf{w}^\top \mathbf{x}(s)$. Let $X \in \R^{\abs{\States} \times m}$ be
  the matrix whose rows are the $\mathbf{x}(s)^\top$ (assumed of full column
  rank) and $D = \mathrm{diag}(d)$. The weighted norm and the \emph{value
  error} are
  \begin{equation}
    \norm{u}_D = \sqrt{\sum_s d(s)\, u(s)^2}, \qquad
    \overline{\mathrm{VE}}(\mathbf{w}) = \norm{X \mathbf{w} - v_\pi}_D^2.
  \end{equation}
  The \emph{projection} $\Pi$ onto the span of the features is $\Pi u = X
  \mathbf{w}_u$ with $\mathbf{w}_u = \argmin_{\mathbf{w}} \norm{X
  \mathbf{w} - u}_D$, i.e.\ $\Pi = X (X^\top D X)^{-1} X^\top D$.
\end{definition}

\begin{intuition}
  With far fewer parameters than states, the approximation cannot be right
  everywhere, so we must say where errors matter: the weights $d(s)$ make
  the frequently visited states count most. The features are all that the
  approximation ``sees'' of a state: two states with the same features get
  the same value. $\Pi u$ is the best that can be done for $u$ with these
  features, the closest point to $u$ in the space they span.
\end{intuition}

A table is the special case of one-hot features ($m = \abs{\States}$, $X =
I$). Classical feature constructions include polynomials, Fourier bases,
radial basis functions and \emph{tile coding} (several overlapping grids,
each state activating one tile per grid); a neural network can be seen as
a feature map that is learnt together with the last linear layer.

\section{Gradient Monte Carlo and semi-gradient TD}

Stochastic gradient descent on $\overline{\mathrm{VE}}$ would need
$v_\pi(S_t)$. Monte Carlo replaces it by the return $G_t$, an unbiased
sample:
\begin{equation}
  \mathbf{w} \gets \mathbf{w} + \alpha \left( G_t - \mathbf{w}^\top
  \mathbf{x}(S_t) \right) \mathbf{x}(S_t),
\end{equation}
a true stochastic gradient step, which converges (under the Robbins--Monro
conditions) to the minimiser of $\overline{\mathrm{VE}}$ when the states
are sampled from $d$. Semi-gradient TD(0) uses the TD target instead,
treating it as a constant:
\begin{equation}
  \label{eq:linear-td}
  \mathbf{w} \gets \mathbf{w} + \alpha\, \delta_t\, \mathbf{x}(S_t),
  \qquad
  \delta_t = R_{t+1} + \gamma\, \mathbf{w}^\top \mathbf{x}(S_{t+1}) -
  \mathbf{w}^\top \mathbf{x}(S_t).
\end{equation}

\begin{intuition}
  The update moves the prediction for $S_t$ towards the target, through the
  features of $S_t$. It is only a ``semi''-gradient because the target also
  depends on $\mathbf{w}$ (through $S_{t+1}$), and that dependence is
  ignored: we move the prediction towards the target, never the target
  towards the prediction. A full gradient would also change the value of
  $S_{t+1}$ to make it agree with $S_t$, letting the future be distorted to
  fit the past, which is not what we want.
\end{intuition}

\begin{proposition}{Expected TD update}{td-expected}
  Let $P_\pi$ and $r_\pi$ be as in \cref{prop:matrix}. If $S_t \sim d$ and
  the transition follows $\pi$, then
  \begin{equation}
    \E\left[\delta_t\, \mathbf{x}(S_t)\right] = \mathbf{b} - A \mathbf{w},
    \qquad A = X^\top D (I - \gamma P_\pi) X, \qquad \mathbf{b} = X^\top D
    r_\pi.
  \end{equation}
\end{proposition}

\begin{proof}
  Conditioning on $S_t = s$,
  \begin{equation*}
    \E[\delta_t \mathbf{x}(S_t)] = \sum_s d(s)\, \mathbf{x}(s) \left(
    r_\pi(s) + \gamma \sum_{s'} P_\pi(s, s')\, \mathbf{x}(s')^\top
    \mathbf{w} - \mathbf{x}(s)^\top \mathbf{w} \right),
  \end{equation*}
  which is the stated matrix expression.
\end{proof}

If TD converges, it is to a point where the expected update vanishes: the
\emph{TD fixed point} $\mathbf{w}_{\mathrm{TD}} = A^{-1} \mathbf{b}$. It is
generally \emph{not} the minimiser of $\overline{\mathrm{VE}}$.

\section{The TD fixed point}

\begin{lemma}{$P_\pi$ is non-expansive in $\norm{\cdot}_D$}{p-nonexpansive}
  If $d P_\pi \leq d$ componentwise (true if $d$ is the stationary
  distribution of $P_\pi$, or the expected number of visits of each state
  in an episodic task), then $\norm{P_\pi u}_D \leq \norm{u}_D$ for every
  $u$.
\end{lemma}

\begin{proof}
  Each row of $P_\pi$ is a (sub-)probability vector, so by Jensen's
  inequality (convexity of $x \mapsto x^2$, and the missing mass counting
  as a value $0$) $\left( \sum_{s'} P_\pi(s, s') u(s') \right)^2 \leq
  \sum_{s'} P_\pi(s, s')\, u(s')^2$. Hence
  \begin{equation*}
    \norm{P_\pi u}_D^2 \leq \sum_s d(s) \sum_{s'} P_\pi(s, s')\, u(s')^2 =
    \sum_{s'} (d P_\pi)(s')\, u(s')^2 \leq \sum_{s'} d(s')\, u(s')^2.
    \qedhere
  \end{equation*}
\end{proof}

\begin{theorem}{Convergence of on-policy linear TD}{linear-td}
  Under the assumption of \cref{lem:p-nonexpansive}, the matrix $A$ is
  positive definite: $\mathbf{w}^\top A \mathbf{w} \geq (1 - \gamma)
  \norm{X \mathbf{w}}_D^2 > 0$ for $\mathbf{w} \neq 0$. Consequently the
  expected update $\mathbf{w}_{k+1} = \mathbf{w}_k + \alpha (\mathbf{b} - A
  \mathbf{w}_k)$ converges to $\mathbf{w}_{\mathrm{TD}}$ for every small
  enough $\alpha > 0$.
\end{theorem}

\begin{proof}
  Let $\mathbf{y} = X \mathbf{w}$. By Cauchy--Schwarz in the inner product
  $\langle u, v \rangle_D = u^\top D v$ and \cref{lem:p-nonexpansive},
  \begin{equation*}
    \mathbf{w}^\top A \mathbf{w} = \norm{\mathbf{y}}_D^2 - \gamma
    \langle \mathbf{y}, P_\pi \mathbf{y} \rangle_D
    \geq \norm{\mathbf{y}}_D^2 - \gamma \norm{\mathbf{y}}_D \norm{P_\pi
    \mathbf{y}}_D \geq (1 - \gamma) \norm{\mathbf{y}}_D^2,
  \end{equation*}
  and $\mathbf{y} \neq 0$ when $\mathbf{w} \neq 0$ ($X$ has full column
  rank and we assume $d > 0$). Let $\lambda > 0$ be the smallest eigenvalue
  of the symmetric part $(A + A^\top)/2$, so that $\mathbf{e}^\top A
  \mathbf{e} \geq \lambda \norm{\mathbf{e}}^2$. The error $\mathbf{e}_k =
  \mathbf{w}_k - \mathbf{w}_{\mathrm{TD}}$ satisfies $\mathbf{e}_{k+1} =
  (I - \alpha A) \mathbf{e}_k$, so
  \begin{equation*}
    \norm{\mathbf{e}_{k+1}}^2 = \norm{\mathbf{e}_k}^2 - 2 \alpha
    \mathbf{e}_k^\top A \mathbf{e}_k + \alpha^2 \norm{A \mathbf{e}_k}^2
    \leq \left( 1 - 2 \alpha \lambda + \alpha^2 \norm{A}^2 \right)
    \norm{\mathbf{e}_k}^2,
  \end{equation*}
  and the factor is less than $1$ for $0 < \alpha < 2 \lambda /
  \norm{A}^2$.
\end{proof}

\begin{intuition}
  Positive definiteness of $A$ means that, on average, the update always
  points towards the TD fixed point, like a ball rolling into a bowl. It
  depends on the states being weighted as the policy visits them: then a
  state whose value is overestimated is also visited as a \emph{next}
  state often enough for the overestimate to be corrected.
\end{intuition}

The stochastic algorithm \eqref{eq:linear-td}, with decreasing step sizes,
converges to $\mathbf{w}_{\mathrm{TD}}$ with probability $1$ (Tsitsiklis and
Van Roy, 1997). How good is the fixed point?

\begin{theorem}{Projected Bellman equation and error bound}{td-bound}
  Under the same assumption, $v_{\mathrm{TD}} = X \mathbf{w}_{\mathrm{TD}}$
  is the unique solution of the \emph{projected Bellman equation} $v = \Pi
  \T^\pi v$, and
  \begin{equation}
    \label{eq:td-bound}
    \norm{v_{\mathrm{TD}} - v_\pi}_D \leq \frac{1}{\sqrt{1 - \gamma^2}}\,
    \norm{\Pi v_\pi - v_\pi}_D.
  \end{equation}
\end{theorem}

\begin{proof}
  $\mathbf{b} - A \mathbf{w} = X^\top D (r_\pi + \gamma P_\pi X \mathbf{w} -
  X \mathbf{w}) = X^\top D (\T^\pi v - v)$ with $v = X \mathbf{w}$. It
  vanishes iff $\T^\pi v - v$ is $D$-orthogonal to the span of $X$, i.e.\
  iff $\Pi(\T^\pi v - v) = 0$, i.e.\ $v = \Pi \T^\pi v$ since $\Pi v = v$.
  The map $\Pi \T^\pi$ is a $\gamma$-contraction in $\norm{\cdot}_D$
  ($\Pi$ is an orthogonal projection, hence non-expansive, and $\norm{\T^\pi
  u - \T^\pi u'}_D = \gamma \norm{P_\pi (u - u')}_D \leq \gamma \norm{u -
  u'}_D$), which gives uniqueness. For the bound, $v_\pi - \Pi v_\pi$ is
  orthogonal to the span, which contains $\Pi v_\pi - v_{\mathrm{TD}}$; by
  Pythagoras and the contraction, using $v_\pi = \T^\pi v_\pi$,
  \begin{align*}
    \norm{v_\pi - v_{\mathrm{TD}}}_D^2
    &= \norm{v_\pi - \Pi v_\pi}_D^2 + \norm{\Pi \T^\pi v_\pi - \Pi \T^\pi
      v_{\mathrm{TD}}}_D^2 \\
    &\leq \norm{v_\pi - \Pi v_\pi}_D^2 + \gamma^2 \norm{v_\pi -
      v_{\mathrm{TD}}}_D^2,
  \end{align*}
  and we solve for $\norm{v_\pi - v_{\mathrm{TD}}}_D^2$.
\end{proof}

\begin{intuition}
  TD never sees $v_\pi$: it only sees its own predictions one step ahead.
  Looking one step ahead ($\T^\pi$) produces a function that the features
  may not represent, so it is projected back ($\Pi$); the TD solution is the
  function that this look-then-project operation leaves unchanged. It is
  generally not the best fit $\Pi v_\pi$, because the approximation errors
  of the next states feed into the targets. The bound says that this
  compounding costs at most a factor $1 / \sqrt{1 - \gamma^2}$ ($7$ for
  $\gamma = 0.99$), and often much less.
\end{intuition}

\begin{pacman}{Linear features for the reduced game}{vfa-escape}
  We approximate the value $v_\star$ of the optimal policy of the reduced
  game of \cref{pac:dp-escape} ($5776$ non-terminal states) with $8$
  features of the board: a constant, the maze distances from Pacman to the
  power cookie, from Pacman to the ghost and from the ghost to the cookie,
  the square of the first, the distance to the ghost capped at $6$, and
  two indicators (ghost within $2$ cells; ghost closer to the cookie than
  Pacman). The weighting $d$ is the expected number of visits from a
  uniform start, which satisfies $d P_\pi \leq d$. The script
  \code{Course/code/vfa\_escape.py} computes exactly:
  \begin{center}
    \begin{tabular}{@{}lc@{}}
      \toprule
      Quantity & Value \\
      \midrule
      $\norm{v_\star}_D$ & $8.261$ \\
      best approximation $\norm{\Pi v_\star - v_\star}_D$ & $0.245$ \\
      TD fixed point $\norm{v_{\mathrm{TD}} - v_\star}_D$ & $0.262$ \\
      bound \eqref{eq:td-bound} & $1.734$ \\
      \bottomrule
    \end{tabular}
  \end{center}
  Eight numbers summarise $5776$ values to within $3\%$, and the TD fixed
  point is almost as good as the best approximation: the bound, which must
  hold for every problem, is pessimistic here. The smallest eigenvalue of
  the symmetric part of $A$ is $0.005$, positive as
  \cref{thm:linear-td} predicts.
\end{pacman}

\section{Least-squares TD}

Since $\mathbf{w}_{\mathrm{TD}} = A^{-1} \mathbf{b}$, one can estimate $A$
and $\mathbf{b}$ from data and solve the linear system directly
(LSTD, Bradtke and Barto, 1996):
\begin{equation}
  \hat A = \sum_t \mathbf{x}(S_t) \left( \mathbf{x}(S_t) - \gamma\,
  \mathbf{x}(S_{t+1}) \right)^\top + \varepsilon I, \qquad
  \hat{\mathbf{b}} = \sum_t R_{t+1}\, \mathbf{x}(S_t), \qquad
  \mathbf{w} = \hat A^{-1} \hat{\mathbf{b}}.
\end{equation}
LSTD has no step size and uses the data optimally, at a cost of $O(m^2)$
per step (with incremental inverses) instead of $O(m)$; it is attractive
for moderate numbers of features.

\section{Off-policy divergence}

The proof of \cref{thm:linear-td} used $d P_\pi \leq d$: the states must be
weighted as the policy being evaluated visits them. With another weighting
(off-policy learning, or a replay memory filled by other policies), $A$ can
fail to be positive definite and TD can diverge. \cref{prop:divergence}
gives a two-state example; the classical one is \emph{Baird's
counterexample} (1995): seven states, linear features with more parameters
than needed, uniform off-policy updates, and the weights diverge for every
positive step size, although an exact representation exists. This is the
\emph{deadly triad} of function approximation, bootstrapping and
off-policy learning.

\begin{intuition}
  In the off-policy case, some states are updated often while the states
  their targets depend on are rarely updated. Through the shared weights,
  raising the value of the first states also raises the value of the
  second, which raises the targets of the first: the updates chase a target
  that they push away. On-policy, every state that serves as a target is
  also updated in proportion, and this feedback is kept in check.
\end{intuition}

Two families of remedies exist. \emph{Gradient TD} methods (GTD2, TDC,
Sutton et al., 2009) perform true stochastic gradient descent on the
\emph{mean squared projected Bellman error}
\begin{equation}
  \overline{\mathrm{PBE}}(\mathbf{w}) = \norm{\Pi \T^\pi X \mathbf{w} - X
  \mathbf{w}}_D^2 = (\mathbf{b} - A \mathbf{w})^\top (X^\top D X)^{-1}
  (\mathbf{b} - A \mathbf{w}),
\end{equation}
estimating the matrix product with a second set of weights; they converge
off-policy. \emph{Emphatic} TD reweights the updates so that the effective
weighting satisfies the condition of \cref{lem:p-nonexpansive}. Deep RL
mostly relies instead on the stabilisers of DQN (target networks, replay
memories close to on-policy), which weaken the triad without removing it.

\section{Control with function approximation}

For control, $\hat q(s, a; \mathbf{w})$ is learnt with the semi-gradient
versions of SARSA or Q-learning. Semi-gradient SARSA with linear features
can oscillate but stays bounded (Gordon, 2001); Q-learning with linear
features can diverge, being off-policy. In practice, the combination of
Q-learning with deep networks requires the techniques of the next chapter.

\section{Summary}

\begin{itemize}[nosep]
  \item Linear approximation represents values as $\mathbf{w}^\top
    \mathbf{x}(s)$; the error is measured in the norm weighted by the
    state distribution.
  \item On-policy, semi-gradient TD converges to the TD fixed point, the
    solution of the projected Bellman equation, within a factor $1 /
    \sqrt{1 - \gamma^2}$ of the best approximation; LSTD computes it
    directly.
  \item Off-policy, TD can diverge (deadly triad); gradient TD methods
    converge, and deep RL relies on target networks and replay memories.
\end{itemize}
